\documentclass{article}
\usepackage[a4paper,landscape,margin=0.8cm]{geometry}
\usepackage[T1]{fontenc}
\usepackage[utf8]{inputenc}
\usepackage[spanish]{babel}
\usepackage{newtxtext,newtxmath}   % Times New Roman
\usepackage{multicol}
\usepackage{xcolor}
\pagestyle{empty}
\setlength{\parindent}{0pt}
\setlength{\parskip}{2pt}
\setlength{\columnsep}{0.8cm}
\setlength{\columnseprule}{0.3pt}
\newcommand{\ej}[1]{\par\smallskip\textbf{Ejercicio #1.}\ }
\newcommand{\rta}{\par\textbf{Respuesta:}\ }
\newcommand{\cod}[1]{{\ttfamily\fontsize{6.5}{7.5}\selectfont #1}}

\begin{document}
\color{black}
\fontsize{8}{9.6}\selectfont
\begin{multicols}{2}

\ej{1} \cod{xTaskCreate(tareaRed,"RED",3072,NULL,3,\&hRed); xTaskCreate(tareaLog,"LOG",2048,NULL,1,\&hLog);}
¿Qué son los 6 argumentos, quién tiene más prioridad y RED se ejecuta siempre?
\rta (1) función de la tarea, (2) nombre, (3) tamaño de pila (en bytes en ESP-IDF/Arduino-ESP32), (4) parámetro, (5) prioridad, (6) puntero donde se guarda el handle. RED (3) tiene mayor prioridad que LOG (1). No se ejecuta siempre: puede bloquearse, suspenderse o ser desplazada por otra de mayor prioridad. Se debe verificar que retorne \cod{pdPASS}.

\ej{2} A (prio 4): \cod{procesar(); vTaskDelay(10 ms);} \ B (prio 2): \cod{registrar(); vTaskDelay(50 ms);}. A acaba de llamar a \cod{vTaskDelay} y B está lista. ¿Quién ejecuta y qué pasa al vencer el retraso de A?
\rta A pasa de Running a Blocked y B pasa de Ready a Running. A los 10 ms A vuelve a Ready y, por tener mayor prioridad, desplaza a B de inmediato. \cod{vTaskDelay} no consume CPU: bloquea la tarea.

\ej{3} \cod{for(;;)\{ trabajo(); vTaskDelay(pdMS\_TO\_TICKS(20)); \}}, el trabajo dura 7 ms. ¿Intervalo entre inicios? ¿Cómo lograr período de 20 ms?
\rta Intervalo $\approx 7+20=27$ ms, porque \cod{vTaskDelay} espera 20 ms \emph{después} del trabajo (hay deriva). Para 20 ms fijos usar \cod{vTaskDelayUntil}:
\cod{TickType\_t ant=xTaskGetTickCount(); for(;;)\{ trabajo(); vTaskDelayUntil(\&ant,pdMS\_TO\_TICKS(20)); \}}. El trabajo debe terminar antes del siguiente instante.

\ej{4} En t=0: A (prio 3) ejecuta 4 ms y bloquea 8 ms (de t=4 a t=12), luego necesita 2 ms más. B (prio 2) necesita 6 ms. C (prio 1) necesita 3 ms. Ejecución entre t=0 y t=13.
\rta $[0,4)$: A. \ $[4,10)$: B (termina). \ $[10,12)$: C (2 de sus 3 ms). \ En t=12 A vuelve a Ready y desplaza a C. \ $[12,13)$: A (1 de sus 2 ms; terminaría en t=14). C queda Ready con 1 ms pendiente. La prioridad solo decide entre tareas listas.

\ej{5} \cod{alta(): while(true)\{ calcular(); \}} (prio 5), \cod{baja(): while(true)\{ atender(); \}} (prio 1), mismo núcleo. ¿Qué ocurre y cómo corregir?
\rta \emph{alta} nunca se bloquea, así que \emph{baja} sufre inanición (nunca obtiene CPU). Poner delay en \emph{baja} no sirve. Se corrige en \emph{alta}: que espere un evento, cola, notificación o use \cod{vTaskDelay} si es periódica; o replantear prioridades. \cod{yield()} no garantiza que corra una tarea de menor prioridad.

\ej{6} A y B (prio 2) listas sin bloqueos; C (prio 3) bloqueada hasta t=5 ms. ¿Cómo se reparte la CPU? ¿Se puede predecir el orden de los \cod{Serial}?
\rta Antes de t=5, A y B se alternan por \emph{time slicing} (en cada tick). En t=5 C se desbloquea y desplaza a A y B por mayor prioridad. No se puede predecir el orden exacto de las impresiones: el puerto serie añade latencia, bloqueos e intercalados.

\columnbreak

\ej{7} \cod{xTaskCreate(tarea,"T",2048,NULL,2,\&h); vTaskSuspend(h); /*100 ms*/ vTaskResume(h);} ¿Transiciones de estado? ¿Los 100 ms la despiertan solos?
\rta Al crearse, T pasa a Ready (y luego puede ir a Running). \cod{vTaskSuspend} la pasa a Suspended; el paso del tiempo \textbf{no} la despierta. \cod{vTaskResume} la devuelve a Ready y correrá cuando el planificador la elija. Suspended $\neq$ Blocked: Blocked sí tiene condición de tiempo o evento para salir.

\ej{8} X (prio 4) espera 100 ms por ciclo; Y (prio 1) calcula 15 ms y espera 5 ms. ¿X consume más CPU por tener mayor prioridad?
\rta No. La prioridad solo define quién va primero cuando ambas están listas, no cuánta CPU recibe. X pasa casi todo el tiempo bloqueada y usa poca CPU; Y aprovecha la CPU mientras X está bloqueada. La carga real se determina midiendo.

\ej{9} \cod{BaseType\_t r=xTaskCreate(tarea,"T",1024,NULL,2,NULL); if(r!=pdPASS) Serial.println("ERROR");} ¿Qué significa ERROR? ¿Siempre faltó memoria? ¿Riesgo de pila pequeña?
\rta ERROR indica que la tarea no se pudo crear. La falta de memoria es la causa habitual, pero no se puede afirmar sin diagnosticar. Una pila pequeña puede \textbf{desbordarse} después, aunque la creación haya sido exitosa. En ESP-IDF la pila se da en bytes; en FreeRTOS estándar, en palabras.

\ej{10} \cod{\#define CONFIG\_TICK\_RATE\_HZ 100} y \cod{vTaskDelay(pdMS\_TO\_TICKS(15));} ¿Cuántos ticks da? ¿Espera exactamente 15 ms?
\rta Tick = 10 ms. $15\cdot100/1000=1.5\to$ \textbf{1 tick} (división entera). No garantiza 15 ms: la espera real depende de la fase del tick (entre 0 y 10 ms aprox.). Para asegurar un mínimo hay que considerar la resolución y el redondeo.

\ej{11} \cod{monitor(): for(;;)\{ Serial.println("OK"); vTaskDelay(pdMS\_TO\_TICKS(100)); \}}. Se miden intervalos de 101, 107, 99 y 112 ms. Tres causas y una mejor medición.
\rta Causas: (1) tiempo de \cod{Serial.println} y su buffer, (2) interferencia de tareas de mayor prioridad, (3) variación del planificador/tick y carga del sistema. Mejor medición: tomar marcas con \cod{esp\_timer\_get\_time()} y calcular diferencias, separar la medición del envío serie, repetir muestras y reportar mínimo, máximo y promedio. Esto no prueba que FreeRTOS falló.

\ej{12} RX (prio 4) recibe tramas cada 10 ms; PROCESA (prio 3) tarda 6 ms por trama; LOG (prio 1) tarda 3 ms por registro. Proponga estados y prioridades, ¿qué pasa si RX sondea sin bloquearse? ¿Qué medir?
\rta Las prioridades son razonables (RX es la más sensible al tiempo). Cada tarea debe estar \textbf{bloqueada} esperando: RX por recepción/evento, PROCESA por mensajes y LOG por registros. Si RX sondea continuamente con prio 4, acapara la CPU y PROCESA y LOG no ejecutan: se acumulan datos y se pierden tramas. La prioridad no reemplaza un diseño con bloqueo. Medir: tasa de recepción, latencia RX$\to$procesamiento, ocupación y desbordamiento de buffers, uso de CPU y variación temporal bajo carga.

\end{multicols}
\end{document}
